\section{Iterative Preorder Traversal}
\label{sec:trees-iter-preorder}

\problemstatement{Return the preorder traversal (node, left, right) of a
binary tree \emph{without} recursion, using an explicit stack.}

\begin{patternbox}
\emph{``Do this traversal iteratively / without recursion''} means:
simulate the call stack yourself with an explicit \textbf{stack}. Preorder
is the easiest of the three to convert because the node is recorded the
instant it is seen, before either child -- so there is no ``come back to
this node later'' bookkeeping to manage.
\end{patternbox}

\subsection*{Understanding the problem carefully}

Recursion works by an implicit stack of pending calls. To remove the
recursion we make that stack explicit. In preorder we want to output a
node, then fully handle its left subtree, then its right subtree. A stack
(LIFO) lets us achieve exactly this if we are careful about push order.

\begin{ideabox}[Push Right Before Left So Left Pops First]
Pop a node and immediately record it. Then push its \emph{right} child and
\emph{then} its \emph{left} child. Because a stack returns the most
recently pushed item first, the left child is processed before the right
-- matching preorder's ``node, then all of left, then all of right.''
Pushing right first is the whole trick; reverse it and you would get
``node, right, left.''
\end{ideabox}

\subsection*{Algorithm}

\begin{enumerate}
  \item If root is null, return empty. Push root onto a stack.
  \item While the stack is non-empty:
  \begin{enumerate}
    \item Pop a node; append its value to the output.
    \item If it has a right child, push it.
    \item If it has a left child, push it.
  \end{enumerate}
\end{enumerate}

\subsection*{Dry run}

\dryrun{Tree: root $1$; children $2$ (left), $3$ (right); $2$'s left child
$4$.}

\begin{itemize}
  \item Stack $=[1]$: pop $1$, record $1$, push $3$ then $2$. Stack
        $=[3,2]$ (top is $2$).
  \item Pop $2$, record $2$, push (no right), push $4$. Stack $=[3,4]$.
  \item Pop $4$, record $4$ (no children). Stack $=[3]$.
  \item Pop $3$, record $3$. Stack empty.
  \item Result: $1,2,4,3$.
\end{itemize}

\subsection*{C++ implementation}

\begin{lstlisting}
#include <bits/stdc++.h>
using namespace std;

struct TreeNode {
    int val;
    TreeNode* left;
    TreeNode* right;
    TreeNode(int v) : val(v), left(nullptr), right(nullptr) {}
};

vector<int> preorderIterative(TreeNode* root) {
    vector<int> out;
    if (root == nullptr) return out;

    stack<TreeNode*> st;
    st.push(root);
    while (!st.empty()) {
        TreeNode* node = st.top();
        st.pop();
        out.push_back(node->val);
        if (node->right) st.push(node->right);   // right first
        if (node->left)  st.push(node->left);    // so left pops first
    }
    return out;
}

int main() {
    TreeNode* root = new TreeNode(1);
    root->left = new TreeNode(2);
    root->right = new TreeNode(3);
    root->left->left = new TreeNode(4);

    for (int x : preorderIterative(root)) cout << x << " ";
    cout << "\n";
    return 0;
}
\end{lstlisting}

\begin{complexitybox}
\textbf{Time Complexity.} $O(n)$ -- each node is pushed and popped exactly
once. \textbf{Space Complexity.} $O(h)$ for the stack in the worst case
(the stack never holds more than one node per level of the current path
plus pending right children).
\end{complexitybox}

\begin{takeawaybox}
Preorder is the friendliest traversal to make iterative because ``record''
happens on first sight. The push-right-then-left order is the reusable
detail. Inorder and postorder are harder precisely because their record
step is delayed until after one or both subtrees are done.
\end{takeawaybox}
